% Appendix D: Paper A, from exam/paper_a.md. The questions only; the model answers are a separate
% appendix, as they are a separate file, so the paper can be sat without them in view.

\chapter{Self-Assessment Paper A}
\label{exa:ch:paper}

\noindent\textbf{Time:} 4 hours. \textbf{Closed book.} A basic calculator is permitted.
\textbf{Total: 100 marks.}

This paper exists to test your own skills and knowledge. It is not a qualification and it gates
nothing. It draws on all ten chapters and on both halves of the system, so \textbf{it is meant to be
taken once you have finished the book}, after \chapref{10} and the bring-up ladder. The model
answers are in Appendix~\ref{ansa:ch:answers}.

\section*{Rubric}
\textbf{The protocol specification is the single source of truth.} Where an implementation and the
protocol disagree, the implementation is wrong. An answer that argues from the contract beats one
that argues from code.

\textbf{Both languages are the subject, so you write both here.} VHDL is VHDL-93 using
\code{ieee.std_logic_1164}, and \code{ieee.numeric_std} where a conversion needs it. C++ that would
run on the ATmega328P is written in the book's AVR-portable subset: \code{<stdint.h>} and the bare
types, no \code{std::}, nested \code{namespace} blocks rather than \code{namespace driver::uart}, no
\code{[[nodiscard]]}, no \code{inline} variables. Host-only test code may use full modern C++, and a
question says so when it means it.

\textbf{Syntax is not what is being marked. The design is.} A missing semicolon, a forgotten
\code{is} or an absent \code{#include} costs nothing. A process that describes a latch where you
meant a wire, or a \code{readReg} that assembles its four bytes backwards, costs that part in full,
because that is a different system.

\textbf{The book's conventions are the ones expected}, and a design that breaks one should say why:
\begin{itemize}
  \item VHDL ports are declared all inputs first, then all outputs, and every \code{port map} and
    \code{generic map} is \textbf{positional}.
  \item An active-low signal is named with an \code{_n} suffix; a signal that has been through a
    two-flop synchronizer is named with an \code{_s2} suffix. Every submodule takes
    \code{reset_s2_n}; only \code{uart_top} takes the raw \code{reset_n}. Signals carrying SPI
    traffic between the two transport blocks take a \code{spi_} prefix.
  \item Resets are asserted \textbf{asynchronously}, so \code{reset_s2_n} belongs in the
    sensitivity list.
  \item C++ uses \code{camelCase}, a leading capital on type names, \code{my} on private members and
    \code{our} on private statics, brace initialization, \code{noexcept} on driver methods, and copy
    and move deleted on any class holding a reference member or a unique hardware resource.
\end{itemize}

\textbf{Where a question asks what reaches the wire, answer in hardware or in bytes}, not in code.
``It calls \code{writeReg}'' is not an answer to ``what does the peripheral do''.

\textbf{Where a question asks for a stated number of defects}, listing more is not penalised, but
only that many are marked, so put your strongest answers first.

\subsection*{Supplied constants and formulas}

\begin{booktable}{@{}p{12em}L@{}}
  \thead{Quantity} & \thead{Value} \\
  \midrule
  DE0-CV system clock & 50\,MHz, a 20\,ns period \\
  Arduino Nano clock (f\textsubscript{osc}) & 16\,MHz \\
  SPI \code{SCK} & f\textsubscript{osc}/16 = 1\,MHz, mode 0, MSB first \\
  Receiver oversampling & 16 times \\
  Default frame & 8N1: 1 start bit, 8 data bits LSB first, 1 stop bit \\
  Baud divider & \code{BAUD_DIV = round( 50_000_000 / (16 * baud) )} \\
  SPI transaction & 5 bytes: 1 command byte, then 4 data bytes MSB first \\
  SPI command byte & bit 7 = \code{W} (1 write, 0 read), bits 6--4 = 0, bits 3--0 = register
    index \\
\end{booktable}

The register map itself is \textbf{not} supplied: Question~\ref{exa:q:1} asks for part of it.

\subsection*{Marks}

\begin{narrowtable}{@{}lcccccccc@{}}
  \thead{Question} & 1 & 2 & 3 & 4 & 5 & 6 & 7 & 8 \\
  \midrule
  \thead{Marks} & 12 & 13 & 14 & 13 & 12 & 13 & 11 & 12 \\
\end{narrowtable}

% ------------------------------------------------------------------------------------------
\question{The map, the top, and what a clean analysis does not prove}{12}
\label{exa:q:1}

\qpart{a} A driver reads \code{STATUS} over SPI and gets back \code{0x0000000D}.

Name every \code{STATUS} bit, give its position, and say whether this value has it set or clear.
Describe the state the peripheral is in, in one sentence. Then state what \code{Uart::write(0x41)}
and \code{Uart::read(byte)} would each do if called immediately afterwards, including any register
access each does \textbf{not} perform.

Finally, give the register index and the byte offset of \code{ERROR_FLAGS}, and the two command
bytes that would read it and write it. \worth{3}

\qpart{b} \code{uart_top}'s entity declares eight ports in the order \code{clock}, \code{reset_n},
\code{sclk}, \code{mosi}, \code{ss}, \code{rx}, \code{miso}, \code{tx}. A colleague swaps
\code{sclk} and \code{mosi} in the entity and changes nothing else.

State whether \code{ghdl -a} rejects the file, and why. State whether elaborating
\code{uart_top_tb} rejects it. Name the first check anywhere in the course that fails, and the
chapter in which that happens.

Then state the single change to the book's conventions that would have moved this failure forward
to analysis time, and what that change costs. \worth{3}

\qpart{c} Two things about \code{reg_rdata} and one about \code{baud_div}.

\medskip\noindent\textbf{(i)} \chapref{1} leaves the placeholder \code{reg_rdata <= (others => '0');}
in \code{uart_top}, and \chapref{5} instantiates \code{uart_regs}, whose \code{reg_rdata} output
drives the same signal. A candidate forgets to delete the placeholder. \code{uart_top_tb} writes
\code{BAUD_DIV} with \code{0x00000004} and reads it back.

Name the VHDL mechanism that now governs \code{reg_rdata}, give the value each of the low four bits
carries during the read-back, and state what the testbench reports. \worth{2}

\medskip\noindent\textbf{(ii)} Before \chapref{5}, \code{baud_div} has no driver at all, and
\code{uart_top} still has to hand \code{baud_gen} a \code{natural range 1 to 65535}. State the value
\code{to_integer(unsigned(baud_div))} returns in that situation and the message it prints, and state
why the value the library picks is worse than merely wrong here. Then give the guarded expression
the book uses, and say which of its two cases covers this and which covers a legitimate value that
arrives once the register bank exists. \worth{2}

\qpart{d} \code{reset_sync} and \code{sync} are both two flip-flops in series, and the book keeps
them as separate modules.

Name the port on \code{sync} that makes it impossible for \code{sync} to be the reset synchronizer,
and describe the circular dependency you would create by trying. Then state what ``assert
asynchronously, release synchronously'' buys in a design of a few hundred flip-flops, and what could
go wrong if the release were asynchronous too. \worth{2}

% ------------------------------------------------------------------------------------------
\question{The wire, the divider, and the transmitter}{13}
\label{exa:q:2}

\qpart{a} The byte \code{0x53} is sent 8N1 by \code{uart_tx}, which builds its frame as

\begin{vhdlcode}
frame <= STOP_BIT & data & START_BIT;
\end{vhdlcode}

\noindent and then drives \code{tx <= frame(bit_idx)} with \code{bit_idx} walking from 0 to 9.

Give \code{frame} as a 10-bit vector, index 9 down to 0. Then give the ten levels \code{tx} carries
in \textbf{time order}, naming the \code{frame} index each one comes from and what each bit is
(start, data bit~\emph{n}, or stop).

State how many \code{baud_tick}s the whole frame takes, and how many 50\,MHz clock cycles that is at
115200~baud. \worth{4}

\qpart{b} The peripheral is to run at 57600~baud from the 50\,MHz clock.

Give the exact (unrounded) divider, the value written to \code{BAUD_DIV}, the baud rate actually
achieved, and the error as a percentage with its sign. State why the receiver at the other end
tolerates it, and name precisely what has to still be true, and at which bit of the frame, for the
byte to be received correctly. \worth{3}

\qpart{c} Here is the state machine of \code{uart_tx}, unchanged:

\begin{vhdlcode}
process(clock, reset_s2_n) is
begin
    if (reset_s2_n = '0') then
        tx <= '1'; done <= '0'; state <= STATE_IDLE;
        frame <= (others => '0'); bit_idx <= 0; ticks <= 0;
    elsif (rising_edge(clock)) then
        done <= '0';
        case (state) is
            when STATE_IDLE =>
                tx <= '1';
                if (start = '1') then
                    state <= STATE_SEND;
                    tx    <= '0';
                    frame <= STOP_BIT & data & START_BIT;
                end if;
            when STATE_SEND =>
                tx <= frame(bit_idx);
                if (baud_tick = '1') then
                    if (ticks >= MAX_TICKS - 1) then
                        ticks <= 0;
                        if (bit_idx >= FRAME_WIDTH-1) then
                            done <= '1'; state <= STATE_IDLE; bit_idx <= 0;
                        else
                            bit_idx <= bit_idx + 1;
                        end if;
                    else
                        ticks <= ticks + 1;
                    end if;
                end if;
            when others =>
                state <= STATE_IDLE;
        end case;
    end if;
end process;
\end{vhdlcode}

\noindent State at which clock edge \code{tx} first goes low relative to the edge on which
\code{start} was seen, and say what that means for when the start bit begins relative to
\code{baud_gen}'s tick.

Then, for a divider \code{div}, give the \textbf{shortest} and the \textbf{longest} the start bit
can be, in clock cycles, and explain where the variation comes from. Express the worst case as a
fraction of a nominal bit period, and state why neither the data bits nor the stop bit inherit it.
\worth{3}

\qpart{d} \code{uart_tx_tb} checks the eight data bits least significant first, and \chapref{2}
describes this as the check that catches a transmitter sending most significant first. An earlier
revision of the testbench sent \code{TXBYTE = x"A5"}; it sends \code{x"53"} today.

Write out the eight data-bit levels the line would carry under each of the two orderings for
\code{0xA5}, and use them to show that the testbench could not have caught that bug at all. Name
the property of \code{0xA5} responsible, and state how many of the 256 possible bytes share it.

Then repeat the exercise for \code{0x53}, state which data bit the testbench reports as wrong first,
and say what the episode tells you about reading a passing test as proof. \worth{3}

% ------------------------------------------------------------------------------------------
\question{The asynchronous line}{14}
\label{exa:q:3}

\qpart{a} The \code{rx} pin arrives from another chip.

State what can happen to a flip-flop that samples it directly, and be precise about what the failure
actually is: getting either level back from a mid-transition sample is \emph{not} it. State why two
flip-flops in series fix it and why one does not.

Then state which module in this design instantiates the synchronizer, which module deliberately
does \textbf{not}, and what the receiver's port name for the line tells a reader about that
decision. \worth{3}

\qpart{b} A correctly built \code{uart_rx} receives \code{0x53}, sent 8N1 least significant first,
and stores each sampled bit with \code{frame(bit_idx) <= rx_s2}, walking \code{bit_idx} from 0
upward.

Give the eight data-bit levels in the order they arrive on the line, and the \code{bit_idx} each
lands in. Show \code{frame} after the eighth, and confirm \code{data_out} carries \code{0x53}.

Then a colleague counts \code{bit_idx} \textbf{down} from 7 instead. Give the byte \code{data_out}
would carry for this same frame, in hex, and state the general relationship between what was sent
and what is delivered. \worth{3}

\qpart{c} The receiver aims to sample each bit at tick~8 of 16. It does not quite. Three things
each cost one tick: the falling edge is only noticed at the next \code{baud_tick}; \code{ticks} is
compared before it is incremented; and every \code{ticks <= 0} on a state change spends one more.
The result, as built, is that the start-bit re-check arrives \textbf{nine} ticks after the detected
edge rather than eight, the first data sample \textbf{seventeen} ticks after that re-check, the
stop-bit sample \textbf{seventeen} after the last data bit, and data bit to data bit exactly
\textbf{sixteen}.

\medskip\noindent\textbf{(i)} State how far past the centre of its own bit each of the start-bit
re-check, a data bit and the stop bit is sampled, in ticks. State why the offset is fixed rather
than accumulating across the byte. \worth{1}

\medskip\noindent\textbf{(ii)} Take an \textbf{ideal} receiver first, sampling every bit exactly at
its centre. Working in ticks measured from the start-bit edge, give the tick at which the stop bit
is sampled and the two ticks at which the stop bit begins and ends. Hence derive the combined baud
mismatch, as a percentage, at which the stop-bit sample first falls outside the stop bit.
\worth{2}

\medskip\noindent\textbf{(iii)} Now do the same arithmetic for the receiver as built. Give the two
limits, state which side of the ideal figure has been eaten into and by how much, and say in one
sentence which physical situation that lost margin makes more dangerous. \worth{2}

\qpart{d} Both FIFOs in the finished peripheral are \code{WIDTH = 8, DEPTH = 8}.

Write \code{fifo}'s clocked process, and the concurrent assignments that go with it for
\code{rdata}, \code{empty} and \code{full}. Assume the declarations \code{entries}, \code{head},
\code{tail}, \code{count} and a helper that advances a pointer with wraparound; write everything
else. Guard both directions --- a write to a full FIFO and a read from an empty one must each do
nothing --- and be careful that a push and a pop landing on the same edge leave \code{count} where
it was.

Then software stops reading while frames keep arriving. State how many back-to-back frames the RX
path absorbs before anything is lost, what happens to the next one, and which line of what you just
wrote decides it. Name the \code{ERROR_FLAGS} bit reserved for that loss, and state what it reads in
this book's build and why that is a reservation rather than an omission. \worth{3}

% ------------------------------------------------------------------------------------------
\question{Ports become registers}{13}
\label{exa:q:4}

\qpart{a} \code{STATUS} is computed, never stored.

Write the concurrent assignments that build the whole 32-bit \code{status_reg} in \code{uart_regs}:
all four defined bits, indexed by their \code{uart_def} names, in terms of the signals the bank
actually has, and the reserved bits above them. Then say in a few words, for each defined bit, why
it could not be a stored register that the datapath writes.

Then: \code{STATUS} is read at a moment when the TX FIFO holds two of its eight bytes, the
transmitter is mid-frame, the RX FIFO is empty and no error has been latched. Give the 32-bit value
the bank presents on \code{reg_rdata}. \worth{3}

\qpart{b} A candidate makes \code{RX_DATA} advance the FIFO as a side effect of being read, so that
the driver needs no \code{RX_POP}.

Name the check in \code{uart_regs_tb} that fails and state what it asserts. Then, in terms of the SPI
transport's abort rule, state why a read with a side effect is genuinely harder to get right than a
pure read plus a separate write, and name the thing the bridge would have to tell the bank that it
currently has no way to say.

Finally trace the driver side, with the popping read in place and the driver's own \code{RX_POP}
write still present. The RX FIFO holds \code{0x41} then \code{0x42}, and nothing more arrives. Give
what \code{Uart::read()} hands back on the first call and on the second, and say what became of
\code{0x42}. \worth{4}

\qpart{c} A candidate gates the TX FIFO push on \code{CTRL} bit~0, reasoning that a disabled
peripheral should not transmit. The design analyzes and elaborates.

State what \code{uart_regs_tb} reports and why. State what \code{uart_top_tb} reports and why, and
say why that second failure message is the less useful of the two.

Then name the fact about \code{uart_regs}' port list that settles, on its own, whether any
\code{CTRL} bit may gate anything at all, and say why the bits exist in \code{uart_def} regardless.
\worth{3}

\qpart{d} The TX feeder in \code{uart_top} is two concurrent lines:

\begin{vhdlcode}
tx_load <= (not tx_empty) and (not tx_busy);
tx_pop  <= tx_load;
\end{vhdlcode}

\noindent State why this loads exactly one byte each time the transmitter goes idle with a non-empty
FIFO, never zero and never two, naming what holds \code{tx_load} low for the rest of the frame.

Then a colleague simplifies it to \code{tx_load <= not tx_empty;}, keeping \code{tx_pop <= tx_load;},
arguing that \code{uart_tx} ignores \code{start} while it is sending anyway. They are right about
\code{start}. Say what happens to a TX FIFO holding four bytes, how many of them reach the wire, and
how long the rest survive. State whether \code{uart_top_tb} catches it, and why. \worth{3}

% ------------------------------------------------------------------------------------------
\question{The driver's contracts}{12}
\label{exa:q:5}

\qpart{a} Both sides of the wire store bit \textbf{positions}, not masks.

Give the C++ expression that tests RX-valid in a \code{uint32_t status}, and the VHDL expression
that indexes the same bit of a \code{std_logic_vector}.

Then three one-line changes. For each, name the defect and state what the running system does ---
not merely that it is wrong.
\begin{itemize}
  \item \code{if (status & status::RX_VALID) { ... }}
  \item \code{writeReg(reg::CTRL, ctrl::ENABLE);}
  \item \file{register_map.hpp} declares \code{constexpr uint8_t RX_VALID{1U << 1};} while
    \file{uart_def.vhd} keeps \code{ST_RX_VALID : natural := 1;}
\end{itemize}
\worth{4}

\qpart{b} \code{driver::transport::Interface} has three methods and knows nothing about registers.

Name the three, and name the two things the seam deliberately does not know. State what putting the
seam at the \textbf{byte} level buys that a seam exposing \code{readReg} and \code{writeReg} would
not, then name the two concrete implementations the book puts under it and state exactly what
changes above the seam when they swap. \worth{3}

\qpart{c} This header compiles on the host and fails on the AVR toolchain the book targets:

\begin{cppcode}
#include <cstdint>

namespace driver::uart
{
inline constexpr std::uint8_t TxReady{0U};

class Interface
{
public:
    virtual ~Interface() noexcept = default;
    [[nodiscard]] virtual std::uint32_t status() const noexcept = 0;
};
} // namespace driver::uart
\end{cppcode}

\noindent Name \textbf{four} things in it that the AVR build rejects, and in each case say what the
toolchain is actually missing rather than merely that it is old.

Then write the header out corrected, as the book would have it: the include, the namespaces, the
constant and the class, with nothing left in it the AVR build cannot take. \worth{3}

\qpart{d} \code{write()} and \code{read()} are non-blocking and each returns a \code{bool}.

State what each \code{bool} means. Name the \code{STATUS} bit that makes each answer possible, and
the \chapref{5} mechanism behind that bit. Then state where the blocking versions live, why they are
free functions taking \code{Interface&} rather than methods, and what would be lost if the interface
itself blocked. \worth{2}

% ------------------------------------------------------------------------------------------
\question{The five-byte transaction}{13}
\label{exa:q:6}

\qpart{a} A host test constructs a \code{driver::transport::Stub}, constructs a real \code{Uart}
over it, and runs some driver calls. Afterwards the stub's recorded MOSI bytes are, in order:

\begin{console}
82 00 00 00 1B
81 00 00 00 01
00 00 00 00 00
04 00 00 00 00
85 00 00 00 01
\end{console}

\noindent with \code{beginCalls()} and \code{endCalls()} both \code{5}.

Decode every command byte: direction and register. Name the driver call that produced each
transaction, and say which transactions belong together as one call.

State what the driver did with the bytes that came back during the third transaction and during the
fourth, being specific about the order in which it assembled them.

Then state what \code{beginCalls()} and \code{endCalls()} prove that the byte log on its own does
not, and describe the wrong behaviour that a byte log alone would fail to distinguish from this one.
\worth{4}

\qpart{b} A candidate sends the four data bytes least significant first in \code{writeReg}, and
assembles the four returned bytes least significant first in \code{readReg}. The two are consistent
with each other, so a stub round trip through this driver still agrees with itself.

\medskip\noindent\textbf{(i)} Give the five bytes that \code{configure(27)}'s first transaction now
puts on the wire, and the 32-bit value the peripheral commits to \code{BAUD_DIV}. State the 16-bit
value \code{baud_div} therefore carries, what \code{uart_top}'s guarded conversion does with it, and
the baud rate the line actually runs at. \worth{2}

\medskip\noindent\textbf{(ii)} The peripheral presents a \code{STATUS} of \code{0x0000000B}. Give
the \code{uint32_t} this driver ends up holding, and state what \code{write()} returns from then on
and what \code{writeBlocking()} therefore does. \worth{1}

\qpart{c} The receive path is poll \code{STATUS}, read \code{RX_DATA}, write \code{RX_POP}.

State what happens if the \code{RX_POP} write is left out, and what \code{app::EchoNode} does on the
bench as a result. State what happens if the pop is issued \textbf{before} the \code{RX_DATA} read.

Then name the host test that catches a driver which pops even when RX-valid is clear, state what
that pop actually does to this peripheral, and say why the answer ``nothing bad'' is not a defence.
\worth{3}

\qpart{d} \code{readReg} is a \code{const} member, and \code{status()} and \code{errorFlags()} are
\code{const} because of it. The transport's \code{begin()}, \code{transfer()} and \code{end()} are
not \code{const}.

State why those three are not \code{const}. State what
\code{const_cast<transport::Interface&>(myTransport)} does, and why the compiler would in fact accept
the calls without it as the class is written today. Name the change to \code{Uart} that would make
the cast genuinely necessary.

Then state the one circumstance under which casting away \code{const} is undefined behaviour, and
say why this is not it. \worth{3}

% ------------------------------------------------------------------------------------------
\question{The real transport, and the real toolchain}{11}
\label{exa:q:7}

\qpart{a} \code{AvrSpi}'s constructor is three statements:

\begin{cppcode}
DDRB  |= (1U << SCK) | (1U << MOSI) | (1U << SS);
PORTB |= (1U << SS);
SPCR   = (1U << SPE) | (1U << MSTR) | (1U << SPR0);
\end{cppcode}

\noindent State what each statement achieves. The transport contract asks for mode 0, MSB first, and
\code{SCK} at 1\,MHz from a 16\,MHz clock: state how this \code{SPCR} value delivers all three,
naming the bits that are \textbf{not} set and what each one's absence means.

State why \code{MISO} appears in none of the three statements. Then state what would be on the
\code{SS} line without the middle statement, and what the FPGA's bridge would make of the first
transaction. \worth{4}

\qpart{b} A colleague writes:

\begin{cppcode}
uint8_t AvrSpi::transfer(const uint8_t byte) noexcept
{
    SPDR = byte;
    return SPDR;
}
\end{cppcode}

\noindent State the byte this returns and where it came from. State how many \code{SCK} periods must
pass before \code{SPDR} holds the byte that arrived on \code{MISO}, and how long that is at this
bench's \code{SCK} rate.

State what the correct spin waits on, and name the side effect that reading \code{SPSR} and then
\code{SPDR} has. Then say what the driver above would see, on a \code{readReg(STATUS)}, with this
\code{transfer} in place. \worth{3}

\qpart{c} \code{SS}/PB2 is configured as an output even though the transport drives the chip select
itself.

State what the SPI peripheral does if \code{SS} is left an input and something pulls it low mid
transaction, naming the \code{SPCR} bit the hardware changes. State what \code{transfer()}'s spin
does in that moment, and give the symptom a bench would show from then on. \worth{2}

\qpart{d} Name two constructs used freely in the host build that do not survive the freestanding
avr-gcc build, and what replaces each.

Then write out \file{avr/env.cpp}'s definitions: every symbol the file has to supply for this book's
code to link on the ATmega, with the right signatures and \code{extern "C"} where it belongs, and
for each name the ordinary C++ construct that emits the reference. Finally state why the file lives
in \file{avr/} rather than \file{source/}. \worth{2}

% ------------------------------------------------------------------------------------------
\question{The bench}{12}
\label{exa:q:8}

\qpart{a} A user types \code{A} (\code{0x41}) into the PC terminal, and \code{app::EchoNode} echoes
it.

Name, in order, every module the byte passes through from the USB-serial adapter until it reaches
the terminal again. For each, name the chapter that built it (or mark it as provided), and say which
of the two planes --- control or data --- it is on at that moment.

Then name the three representations the byte takes across that journey, and the two places where it
changes from one to the next. \worth{4}

\qpart{b} A bench climbs the ladder. Rungs a, b and c pass. Rung d fails: bytes the driver writes
appear in the terminal as wrong characters, and characters typed produce framing errors rather than
bytes.

State what the ladder has already cleared, rung by rung, and name what is left that rung d is the
first to exercise.

Name the most likely cause, and justify it by saying what rung c does that hides it. Then give the
arithmetic that would confirm or exonerate \code{BAUD_DIV = 27} as the culprit, and say what
magnitude of error you would have to see before the divider itself is to blame. \worth{3}

\qpart{c} Write the \code{main()} that runs on the ATmega328P on bring-up day, in \file{fw/avr/}:
the includes it needs from this book's headers, the objects it constructs and in what order, the
\code{configure()} call for 115200~8N1, the flag it passes, and the call that starts the
application. It allocates nothing and it never returns.

Then state why the flag this \code{main} passes is left \code{false} where the host test's is not,
and why \code{EchoNode::run()} must poll the non-blocking \code{read()} rather than call
\code{readBlocking()} --- connecting that to how the host test manages to end the loop at all.
\worth{3}

\qpart{d} One of the four SPI lines is wired 5\,V straight to a 3.3\,V FPGA pin, with no level
shifter.

State the likely symptom and the risk. Name the first rung of the ladder that exposes it and state
why the rung before it cannot.

Then name the one SPI line for which the omission is electrically safe, say which direction it
travels and why, and state what could still go wrong with it. \worth{2}
